\chapter{Commutativity}
\label{chap:commutativity}


\section{Commutativity of the point measurements}

\begin{theorem}[Commutativity of the point measurements]\label{thm:commutativity-points}
  \lean{MIPStarRE.LDT.CommutativityPoints.commutativityPoints}
  \uses{lem:good-strategy-characterization, prop:simeq-to-approx}
  Let $(\psi,A,B,L)$ be an $(\eps,\delta,\gamma)$-good symmetric strategy for the $(m,q,d)$ low individual degree test. On average over independent $u,v \sim \F_q^m$,
  \[
    (A_a^u A_b^v) \ot I \approx_{32\gamma m} (A_b^v A_a^u) \ot I.
  \]
\end{theorem}
\begin{proof}
  Use the diagonal-line test restricted to the full set of coordinates to compare each point measurement with the common diagonal-line answer. After converting those consistencies to state-dependent distance, commute through the shared diagonal-line measurement and transfer back.
\end{proof}

\section{Commutativity of the slice measurements}

\begin{lemma}[Commutativity after evaluation]\label{lem:comm-data-processed-g}
  \lean{MIPStarRE.LDT.Commutativity.commDataProcessedG}
  \uses{prop:cons-sub-meas, prop:closeness-of-ip, prop:two-notions-of-self-consistency-after-evaluation, thm:commutativity-points}
  Let $\{G^x\}_{x \in \F_q}$ be projective submeasurements in $\polysub{m}{q}{d}$ such that:
  \begin{enumerate}
    \item on average over $(u,x) \sim \F_q^{m+1}$,
    \[
      A_a^{u,x} \ot I \simeq_\zeta I \ot G^x_{[g(u)=a]};
    \]
    \item on average over $x \sim \F_q$, one has $G_g^x \ot I \approx_\zeta I \ot G_g^x$;
    \item there are positive semidefinite operators $Z^x$ with
    \[
      \E_x \bra{\psi} (I-G^x) \ot Z^x \ket{\psi} \le \zeta,
      \qquad
      Z^x \ge \E_u A_{g(u)}^{u,x}.
    \]
  \end{enumerate}
  Write $G_a^{u,x}=G_{[g(u)=a]}^x$. Then, on average over independent $(u,x),(v,y) \sim \F_q^{m+1}$,
  \[
    G_a^{u,x} G_b^{v,y} \ot I \approx_{48m(\gamma^{1/2}+\zeta^{1/2})} G_b^{v,y} G_a^{u,x} \ot I.
  \]
\end{lemma}
\begin{proof}
  Write $G_a^{u,x}=G^x_{[g(u)=a]}$ and expand the commutator norm into the two quartic terms that appear in the source. The bridge to point measurements is not a single substitution: first use consistency to insert the evaluated point projector $A_a^{u,x}$, then use two boundedness-driven stability claims to show that inserting or deleting the total projector $G^x$ or $G^y$ only costs $O(\zeta^{1/2})$. After those internal stability steps, Theorem~\ref{thm:commutativity-points} commutes the point measurements, and the post-processed self-consistency of the $G^x$'s transfers the estimate back to the evaluated families $G_a^{u,x}$.
\end{proof}

\begin{lemma}[Normalization condition for a sandwiched family]\label{lem:normalization-condition}
  \lean{MIPStarRE.LDT.Commutativity.normalizationCondition}
  \uses{def:submeasurement}
  Let $P=\{P_a\}$ be a submeasurement and $Q=\{Q_b\}$ a projective submeasurement. Define $C_{a,b}=Q_b P_a Q_b$. Then
  \[
    \sum_a \left(\sum_b C_{a,b}\right)^\dagger \left(\sum_b C_{a,b}\right) = \sum_a \left(\sum_b C_{a,b}\right) \left(\sum_b C_{a,b}\right)^\dagger
    \le I.
  \]
\end{lemma}
\begin{proof}
  Expand the squares, use the orthogonality of the projectors $Q_b$, and sum over the submeasurement $P$.
\end{proof}

\begin{theorem}[Commutativity of the slice measurements]\label{thm:com-main}
  \lean{MIPStarRE.LDT.Commutativity.comMain}
  \uses{lem:comm-data-processed-g, lem:normalization-condition, prop:switch-sandwich, prop:closeness-of-ip, lem:schwartz-zippel-individual}
  Let $\{G^x\}_{x \in \F_q}$ satisfy the same hypotheses as in Lemma~\ref{lem:comm-data-processed-g}. Then, on average over independent $x,y \sim \F_q$,
  \[
    G_g^x G_h^y \ot I
    \approx_{30m(\gamma^{1/4}+\zeta^{1/4}+(d/q)^{1/4})}
    G_h^y G_g^x \ot I.
  \]
\end{theorem}
\begin{proof}
  Expand the commutator norm into two quartic terms. The first is reduced to the mixed overlap $\bra{\psi} G\ot G\ket{\psi}$ by Lemma~\ref{prop:switch-sandwich}. For the second, the source passes to random evaluations $g(u)$ and $h(v)$, so that Lemma~\ref{lem:comm-data-processed-g} can commute the processed operators. Each application of the inner-product comparison lemma requires the normalization condition from Lemma~\ref{lem:normalization-condition} for sandwiched families such as $G_h^y G_{[g(u)=a]}^x G_h^y$. Finally one removes the two point evaluations by Schwartz--Zippel, paying one $md/q$ loss for each removal. This evaluated-to-polynomial lift is the main work in the theorem.
\end{proof}
